\scanpage{26}
\uline{\unsure{Th.|Pr.}} Soient $S$ un schéma de car. $0$, $X/S$ un schéma, $\widehat{\mathcal P}=\varprojlim\mathcal P_n$ une \deleted{algèbre} $\mathcal O_X$-Algèbre $\mathcal O_X$-augmentée adique complète, \deleted{\unsure{une}} ($\mathcal J$-adique, \inserted{à filtration adique,} donc $\mathcal J/\mathcal J^2$ \deleted{$\unsure{\simeq\Omega}$} est \uline{inversible}), et muni d'une \deleted{$\unsure{\xi}$} connexion \inserted{\unsure{\deleted{très régulière}|\uline{très régulière}}} $C$ rel. $/S$, respectant la structure d'Algèbre, mais \unsure{non néc.} l'augmentation.
\figure{\begin{tikzpicture}\draw (0,0)--(0,-.8);\end{tikzpicture}}
\inserted{On suppose que l'hom \uuline{lin.} associé
\diagram{\begin{tikzcd}
\mathcal J/\mathcal J^2=\omega \arrow[r,"\bar\omega"] & \Omega^1_{X/S}=\unsure{\widehat S\otimes\Omega^1_{X/S}/\mathcal J\widehat S\otimes\Omega^1_{X/S}}
\end{tikzcd}}
\deleted{\unsure{est}} est un iso. \deleted{(\unsure{conn.} \illegible{3})}}

On cherche les fibrés en droites $P$ sur $X$, munis d'une sect. $\sigma$, (donc défini par extension de \unsure{Mod. loc. lib}
\diagram{\begin{tikzcd}[column sep=small]
0 \arrow[r,dash] & \Omega \arrow[r,dash] & E \arrow[r,dash] & \mathcal O \arrow[r,dash] & 0
\end{tikzcd}}
) muni d'une connexion $c$ (ne respectant pas \unsure{néces.} la section \unsure{marquée|canonique} $\sigma$), \deleted{\unsure{tels}} et les iso d'algèbres augm.
\diagram{\begin{tikzcd}
\widehat S \arrow[r,leftrightarrow,"\substack{u\\\sim}"] & \mathcal P
\end{tikzcd}}
\deleted{tel que} (où $\widehat S=\mathcal O_{\widehat P}$, $\widehat P$ complété formel le long de $\sigma$) tels que $u(\unsure{\widehat C|C})=\widehat c$.

a) La donnée d'un iso $u$ à l'ordre $2$ détermine $(P,\sigma)$, i.e. $\Omega,E$, à isom unique près par
\[
\Omega=\mathcal J/\mathcal J^2,\qquad E=\Omega^{-1}\otimes\mathcal J/\mathcal J^3.
\]
b) $(P,\sigma)$ étant ainsi déterminé, pour tt isom \unsure{alg} à l'ordre $3$ près
\diagram{\begin{tikzcd}
\widehat S/\widehat{\mathcal J}_P^{4} \arrow[r] & \mathcal P/\mathcal J^4
\end{tikzcd}}
prolongeant l'iso à l'ordre $2$ canonique, il existe \uline{un isomorphisme et un seul} $u$ qui le prolonge, et tel que $\unsure{\widehat u|u}(C)$ soit de la forme $\widehat c$, et pour \unsure{tte}
% archive foliation: 25
